\chapter{Methodology}
\label{ch:methodology}

\section{Research design}
\textit{Describe the mixed-methods multiple-case design and justify the combination of quantitative analysis and qualitative evidence.}

\section{Case selection}
\textit{Document the prescan protocol, the exploratory feasibility criteria, and the final case lock.}

\section{Phase 1: AI contribution detection}
\textit{Explain the PR-first detection pipeline, AIDev as the primary source, the supplementary marker-based pass, confidence tiers, and manual validation.}

\section{Phase 2: quantitative code and security analysis}
\subsection{Technical debt and code quality metrics}
\textit{Define the metrics, the pre/post-AI windows, the project-level and contribution-level comparisons, and the planned robustness checks.}

\subsection{Security tracing}
\textit{Explain how advisories, fix commits, pull requests, and candidate introducing commits will be linked and graded for confidence.}

\section{Phase 3: qualitative evidence}
\subsection{Public discourse mining}
\textit{Explain how issues, pull requests, release notes, and related discussions will be sampled and coded.}

\subsection{Developer interviews}
\textit{Use this section only if interviews are actually conducted. Otherwise state clearly that public discourse is the primary qualitative source.}

\section{Validity and limitations}
\textit{Address internal validity, construct validity, external validity, measurement error, and the limits of public AI-signal detection.}
